\minitoc 

% ========================================================================================================
% Analyse de Fourier des signaux à temps discret

\section{Transformée de Fourier des signaux à temps discret (TFTD)}

\subsection{Définitions et conventions}

\begin{definition}[TFTD en convention fréquence]
    Pour un signal $x(n)$ à temps discret, sa Transformée de Fourier à Temps Discret est la fonction continue $X(f)$ définie par : 
        \begin{equation}
            \boxed{X(f) = \sum_{n=-\infty}^{+\infty} x(n) e^{-2j \pi n f}}
        \end{equation}
    \textbf{Remarques clés :}
    \begin{itemize}
        \item Bien que $x(n)$ soit discret, $X(f)$ est une fonction continue de la fréquence $f$.
        \item $X(f) \in \C$, avec $X(f) = |X(f)| e^{j \text{Arg}(X(f))}$.
        \item $X(f)$ est périodique de période $1$ : $X(f + 1) = X(f)$.
    \end{itemize}
\end{definition}

\begin{definition}[TFTD inverse]
    La reconstruction du signal temporel s'effectue sur une période fréquentielle unitaire :
        \begin{equation}
            x(n) = \int_{-1/2}^{+1/2} X(f) e^{2j \pi n f} \, df
        \end{equation}
    On note symboliquement cette correspondance : $x(n) \overset{\mathcal{F}}{\longleftrightarrow} X(f)$.
\end{definition}

\begin{remark}[Convention en pulsation $\omega$]
    En posant la pulsation normalisée $\omega = 2\pi f$ ($\omega \in [-\pi, \pi]$) :
    \begin{equation}
        X(\omega) = \sum_{n=-\infty}^{+\infty} x(n) e^{-j n \omega}, \qquad x(n) = \frac{1}{2\pi} \int_{-\pi}^{+\pi} X(\omega) e^{j n \omega} \, d\omega
    \end{equation}
    $X(\omega)$ est alors périodique de période $2\pi$. Le passage s'opère directement par $X(f) = X(\omega)|_{\omega = 2\pi f}$.
\end{remark}

\subsection{Propriétés de la TFTD}

\begin{prop}[Propriétés fondamentales]
    Soient $x(n) \overset{\mathcal{F}}{\longleftrightarrow} X(f)$ et $y(n) \overset{\mathcal{F}}{\longleftrightarrow} Y(f)$ :
    \begin{itemize}
        \item \textbf{Linéarité} : $a x(n) + b y(n) \overset{\mathcal{F}}{\longleftrightarrow} a X(f) + b Y(f)$.
        \item \textbf{Retournement temporel} : $x(-n) \overset{\mathcal{F}}{\longleftrightarrow} X(-f)$.
        \item \textbf{Décalage temporel} : $x(n - n_0) \overset{\mathcal{F}}{\longleftrightarrow} X(f) e^{-2j \pi n_0 f}$.
        \item \textbf{Modulation fréquentielle} : $x(n) e^{2j \pi n f_0} \overset{\mathcal{F}}{\longleftrightarrow} X(f - f_0)$.
        \item \textbf{Convolution temporelle} : $(x * y)(n) \overset{\mathcal{F}}{\longleftrightarrow} X(f) \cdot Y(f)$.
        \item \textbf{Multiplication temporelle} : $x(n) \cdot y(n) \overset{\mathcal{F}}{\longleftrightarrow} (X * Y)(f) = \int_{-1/2}^{1/2} X(u) Y(f - u) \, du$.
        \item \textbf{Théorème de Parseval} : $\sum_{n=-\infty}^{+\infty} |x(n)|^2 = \int_{-1/2}^{1/2} |X(f)|^2 \, df$.
    \end{itemize}
\end{prop}

\begin{proposition}[Symétries de la TFTD]
    \begin{center}
        \renewcommand{\arraystretch}{1.3}
        \begin{tabular}{ll}
            \toprule
            \textbf{Signal temporel} $x(n)$ & \textbf{Spectre} $X(f)$ \\
            \midrule
            Réel & Hermitien : $X(-f) = X^*(f)$ ($|X(f)|$ pair, phase impaire) \\
            Réel et pair & Réel et pair \\
            Réel et impair & Imaginaire pur et impair \\
            Imaginaire pur & Anti-hermitien \\
            Imaginaire et pair & Imaginaire pur et pair \\
            Imaginaire et impair & Réel et impair \\
            \bottomrule
        \end{tabular}
    \end{center}
\end{proposition}

\begin{proposition}[Paires usuelles de transformées]
    \begin{center}
        \renewcommand{\arraystretch}{1.6}
        \begin{tabular}{cc}
            \toprule
            $x(n)$ & $X(f)$ \\
            \midrule
            $\delta(n)$ & $1$ \\
            $\delta(n - n_0)$ & $e^{-2j\pi n_0 f}$ \\
            $1$ & $\delta(f)$ sur $[-1/2, 1/2]$ \\
            $e^{2j \pi f_0 n}$ & $\delta(f - f_0)$ \\
            $\cos(2\pi f_0 n)$ & $\frac{1}{2} [\delta(f - f_0) + \delta(f + f_0)]$ \\
            $\sin(2\pi f_0 n)$ & $\frac{1}{2j} [\delta(f - f_0) - \delta(f + f_0)]$ \\
            $a^n u(n) \quad (|a| < 1)$ & $\displaystyle \frac{1}{1 - a e^{-2j \pi f}}$ \\
            \bottomrule
        \end{tabular}
    \end{center}
\end{proposition}

% ========================================================================================================
% Transformée de Fourier Discrète (TFD)

\section{Transformée de Fourier Discrète (TFD) et Algorithme FFT}

\subsection{Définition de la TFD}

On traite désormais des signaux réels ou complexes de durée finie $N$ indexés sur $\llbracket 0, N-1 \rrbracket$.

\begin{definition}[TFD et TFD inverse]
    La Transformée de Fourier Discrète (TFD) sur $N$ points est donnée par :
        \begin{equation}
            \boxed{X(m) = \sum_{n=0}^{N-1} x(n) e^{-2j \pi \frac{n m}{N}}, \quad m \in \llbracket 0, N-1 \rrbracket}
        \end{equation}
    Sa transformée inverse (ITFD) est définie par :
        \begin{equation}
            \boxed{x(n) = \frac{1}{N} \sum_{m=0}^{N-1} X(m) e^{2j \pi \frac{n m}{N}}, \quad n \in \llbracket 0, N-1 \rrbracket}
        \end{equation}
    On note : $x(n) \overset{\text{TFD}}{\longleftrightarrow} X(m)$.
\end{definition}

\begin{remark}[Échantillonnage et Résolution]
    La TFD correspond à l'échantillonnage de la TFTD continue $X(f)$ aux fréquences discrètes normalisées :
    $$ f_m = \frac{m}{N}, \quad m \in \llbracket 0, N-1 \rrbracket $$
    La \textbf{résolution fréquentielle} d'une TFD à $N$ points est donc :
    \begin{equation}
        \Delta f = \frac{1}{N} \quad (\text{ou } \Delta F = \frac{F_e}{N} \text{ en Hertz})
    \end{equation}
\end{remark}

\begin{prop}[Propriétés spécifiques de la TFD]
    Pour des séquences de longueur $N$ :
    \begin{itemize}
        \item \textbf{Décalage circulaire temporel} : $x((n - n_0) \bmod N) \overset{\text{TFD}}{\longleftrightarrow} X(m) e^{-2j \pi \frac{n_0 m}{N}}$.
        \item \textbf{Modulation} : $x(n) e^{2j \pi \frac{m_0 n}{N}} \overset{\text{TFD}}{\longleftrightarrow} X((m - m_0) \bmod N)$.
        \item \textbf{Convolution circulaire} : $(x \circ y)(n) \overset{\text{TFD}}{\longleftrightarrow} X(m) \cdot Y(m)$.
        \item \textbf{Produit temporel} : $x(n) \cdot y(n) \overset{\text{TFD}}{\longleftrightarrow} \frac{1}{N} (X \circ Y)(m)$.
    \end{itemize}
\end{prop}

\subsection{Algorithme de Fourier Rapide (FFT)}

\begin{proposition}[Gain de complexité]
    Le calcul direct d'une TFD requiert $N^2$ multiplications et additions complexes. 
    L'algorithme de la \textbf{FFT} (Fast Fourier Transform, e.g. Cooley-Tukey) exploite la décomposition par division de l'horizon pour $N = 2^p$ ($p \in \N$) et abaisse le coût algorithmique à :
    \begin{equation}
        \mathcal{O}(N \log_2 N) \quad \text{opérations complexes}
    \end{equation}
\end{proposition}

% ========================================================================================================
% Fenêtrage et apodisation

\section{Troncature spectrale et Fenêtres d'apodisation}

Traiter un signal numérique sur un horizon fini $N$ revient à multiplier le signal théorique infini par une fenêtre temporelle de pondération $w(n)$ :
$$ x_{\text{tronqué}}(n) = x(n) \cdot w(n) \overset{\mathcal{F}}{\longleftrightarrow} X(f) * W(f) $$
Ce produit de convolution par la réponse fréquentielle de la fenêtre $W(f)$ provoque un étalement de l'énergie :
\begin{itemize}
    \item \textbf{Largeur du lobe principal} : régit la résolution fréquentielle (capacité à séparer deux raies proches).
    \item \textbf{Niveau des lobes secondaires} : régit la fuite spectrale (\emph{spectral leakage}) vers les fréquences voisines.
\end{itemize}

\begin{proposition}[Comparatif des fenêtres usuelles d'apodisation sur $N$ points]
    \begin{center}
        \renewcommand{\arraystretch}{1.5}
        \begin{tabular}{llcc}
            \toprule
            \textbf{Fenêtre} & \textbf{Définition temporelle} $w(n), n \in \llbracket 0, N-1 \rrbracket$ & \textbf{Lobe principal} & \textbf{Atténuation} \\
            \midrule
            Rectangulaire & $w(n) = 1$ & $\frac{2}{N}$ & $-13\text{ dB}$ \\
            Triangulaire & $w(n) = 1 - \frac{|n - (N-1)/2|}{(N-1)/2}$ & $\frac{4}{N}$ & $-25\text{ dB}$ \\
            Hamming & $w(n) = 0.54 - 0.46 \cos\left(\frac{2\pi n}{N-1}\right)$ & $\frac{4}{N}$ & $-41\text{ dB}$ \\
            Blackman & $w(n) = 0.42 - 0.5 \cos\left(\frac{2\pi n}{N-1}\right) + 0.08 \cos\left(\frac{4\pi n}{N-1}\right)$ & $\frac{6}{N}$ & $-57\text{ dB}$ \\
            \bottomrule
        \end{tabular}
    \end{center}
\end{proposition}